\section{A finite-volume certificate and the matrix state}\label{sec:state}

The construction will produce a certificate for a damped resolvent of
the finite Anderson matrix with the true potentials. We state that target first, then describe
the numerical state and the information retained while the potential is
revealed. Each finite-volume construction fixes the energy and coupling;
its event may depend on both, while its probability estimate is uniform
in the energy interval. Throughout the proof $d\ge3$ is a fixed integer. Constants
denoted by $c,C$, including constants implicit in $O(\cdot)$ and
$\asymp$, may depend on $d$ and the fixed interval $I_d$. No assertion is uniform as $d\to\infty$.

\subsection{The certificate}

Fix the interval $I_d$ from \eqref{eq:energy-windows},
a coupling $0<\lambda<1$, and integers
\begin{equation}\label{eq:scales}
 \begin{aligned}
 n_0&=\lfloor\log_4(3/\lambda^2)\rfloor,\qquad m\ge4n_0,\\
 s_n&=\frac{\lambda^2}{3}4^{n_0-n}\quad(n_0\le n\le m),\\
 \eta&=s_m,\qquad t_n=s_n+\eta.
 \end{aligned}
\end{equation}
Equivalently, if $u=(\lambda^2/3)4^{n_0}\in(1/4,1]$, then
$s_n=u4^{-n}$ and $\eta=u4^{-m}$. Thus
$\frac14 4^{-n}<s_n\le4^{-n}$ and $t_n\asymp4^{-n}$, with
absolute comparison constants. Let
\[
 L=4^{20m},\qquad \mathbb T_L=(\ZZ/L\ZZ)^d,\qquad N=L^d,
\]
and let $H_m$ be the periodic version of~\eqref{eq:model}, in the
negative-hopping convention:
\[
 (H_m f)(i)=-\sum_{|e|_1=1}f(i+e)+\lambda V_i f(i),
 \qquad V_i\ \text{independent and uniform on }[-1,1].
\]
There is an independent potential at every point of this
$d$-dimensional torus.
For $A\subset\mathbb T_L$, write $P_A$ for coordinate restriction or
projection, as indicated by its domain, and let $\diag_A r$ be the
full matrix with diagonal $r_i$ on $A$ and zero elsewhere.

\begin{proposition}[Finite-volume certificate]\label{prop:certificate}
For each fixed $d\ge3$ there exist $\lambda_0(d)>0$ and constants $c,C>0$
with the following property. For every fixed $0<\lambda<\lambda_0(d)$,
every $E\in I_d$, and every $m\ge4n_0$, one
can use auxiliary independent randomness to construct an event
$\mathcal S_m$, depending on the fixed $E,\lambda$, such that
$\PP(\mathcal S_m)\to1$ as $m\to\infty$.
On $\mathcal S_m$ there are a set $A\subset\mathbb T_L$, $|A|\ge N/2$,
and complex numbers $r_i$, $i\in A$, satisfying
\[
 |r_i|\le C,\qquad c\le\Im r_i\le C,
\]
such that
\[
 \widehat G=(H_m-E-\eta\diag_A r)^{-1}
 \quad\text{exists},\qquad
 \Im\widehat G_{ii}\ge c\quad(i\in A).
\]
For fixed $d$ the constants are independent of $E,\lambda,m$, and
the convergence in probability is uniform in $E\in I_d$ for each
fixed $\lambda$. The failure estimates are also uniformly small for
all $m\ge4n_0$ and $u\in(1/4,1]$ once $n_0$ is sufficiently large.
The potential in $H_m$ consists entirely of the original true variables.
\end{proposition}

The proof occupies Sections~\ref{sec:initialization}--\ref{sec:history}.
Section~\ref{sec:spectral} proves that this certificate implies
Theorem~\ref{thm:main}. The true potential vector retains its original
product law as an unconditional marginal of the enlarged probability
space. The success event may depend on that vector. The spectral proof
uses a feasibility event in the true variables and unconditional averaging.

\subsection{True digits, site states, and retained information}

A uniform random variable $\lambda V_i$ has the representation
\begin{equation}
 \lambda V_i=\sum_{k\ge1}\lambda2^{-k}\varepsilon_{i,k},
 \qquad \PP(\varepsilon_{i,k}=1)
       =\PP(\varepsilon_{i,k}=-1)=\tfrac12,
 \label{eq:history-binary}
\end{equation}
with independent signs over all sites and indices. Resolve the possible
binary ambiguity on its null set by a fixed convention. After the first
$h$ signs at a site have been exposed, and before any further information
about that site has been queried, its remaining sum is uniform on
$[-\lambda2^{-h},\lambda2^{-h}]$, with mean zero and variance
$\lambda^2 4^{-h}/3$. In particular, after $n-n_0$ digits its variance
is $s_n$, and its next digit has amplitude $\sqrt{3s_n/4}$.

At each committed state the physical sites are partitioned into an
active set $A$ and a terminal set $T$. Active sites retain artificial
damping. A committed terminal site has its exact true potential, has
been Schur-eliminated, and never changes again. During the one-digit
stages an active site's real potential value is partial. At the final
endpoint, the active sites carry complete true potentials while retaining
damping.

A subset of the active sites is marked \emph{pending}; a nonpending
active site is \emph{coupled}. Before the final residual sweep, the
construction maintains the following implication: at a coupled site the
installed partial value is exactly the true revealed prefix, and no digit
of its remaining suffix has been queried. At a stage boundary this prefix
has $n-n_0$ digits; during a one-digit sweep it has one more digit at sites
already processed. The final sweep installs the entire remaining true
suffix. A pending site's partial value may contain auxiliary digits, and
the construction makes no freshness claim for its true suffix. The mark
is conservative: the marking rule can cover an active site whose own next
digit was not queried. A pending site can leave this class only by becoming
terminal.

The numerical matrices and the information already revealed are separate
records. Inside a trial, newly terminal physical sites and other matrix
changes are provisional until the trial commits. A failure restores the
pretrial potential values and partition into active and terminal sites,
together with the corresponding matrices, and then applies the trial's
marking rule for pending sites.
All answers already obtained remain in the retained history. Previously
committed terminal sites remain permanent.

\begin{figure}[H]
 \centering
 \begin{tikzpicture}[
  state/.style={draw=black!65,rounded corners=2pt,minimum width=37mm,
                minimum height=16mm,align=center,fill=black!3},
  lab/.style={font=\small,align=center,fill=white,inner sep=2pt},
  >=Stealth,thick]
\node[state] (coupled) at (0,1.35) {Coupled active\\\small true digits only};
\node[state] (pending) at (0,-1.35) {Pending active\\\small may use auxiliary digits};
\node[state,fill=black!9] (terminal) at (7.8,0) {Committed terminal\\\small complete true potential\\\small permanent};
\draw[->] (coupled) -- node[lab,right]{failed trial} (pending);
\draw[->] (pending.east) -- node[lab,below,sloped]{successful clearing} (terminal.south west);
\draw[->] (coupled.east) -- node[lab,above,sloped]{included in a\\successful clearing} (terminal.north west);
\draw[densely dashed,black!40] (-2.35,-2.30) rectangle (2.35,2.30);
\node[font=\small,anchor=south] at (0,2.32) {Active sites retain artificial damping};
\end{tikzpicture}

 \caption{The arrows show possible changes between committed site states.
 A pending mark is conservative: the marked site's own true suffix may
 remain unqueried. The final suffix sweep is reached only when no site is
 pending; after it succeeds, the active sites have complete true potentials
 and retain artificial damping.}
 \label{fig:site-states}
\end{figure}

Initially every site is active and coupled, and every partial value is
zero. The retained-information lemma below supplies the conditional law
of an unqueried input. Sections~\ref{sec:local-sweeps} and
\ref{sec:assembly} verify that the inputs selected by ordinary trials
and clearings are eligible for that law. In the matrix estimates,
``fresh'' means a centered variable with the stated law conditional on
the history available before the chosen update. It is an explicit input
until the corresponding eligibility check has been made.

\paragraph{Retained information.}
The construction uses a separate countable array of independent auxiliary
fair signs. It also uses independent uniform orders of fixed finite
sets: a uniform order of a set $S$ of size $M$ is a uniformly chosen
bijection from $\{1,\ldots,M\}$ to $S$. All these sources are independent
of the true signs and of one another. Only an auxiliary sign that is used,
or the next entry of an order being used, is exposed. Section~\ref{sec:local-sweeps}
specifies the source indices for each ordinary trial.

Number the successive queries and completed measurable decisions by
$q=1,2,\ldots$, and let
$\mathcal F_q^{\mathrm{ret}}$ be the information retained after operation $q$.
Initially $\mathcal F_0^{\mathrm{ret}}$ contains only the fixed parameters and deterministic
initial data. A completed measurable numerical choice counts as one
operation even when it is defined by a limit, as with the later ridge choice
or diagonal reset. Before a query, its type and source index are chosen
measurably from the current history. A query may expose one true digit,
one unused auxiliary sign, the remaining true suffix at a selected site,
or the next entry of a selected finite order that has an unexposed entry.
A whole true suffix is one
operation; under the fixed binary convention, its value determines all
the digits it contains. The history records the queried indices and
answers and every exposed order entry. Numerical updates, choices, and
failure decisions must be measurable from the information recorded by
that time, so adjoining them gives no further information.

The current potential values, active and terminal sets, pending marks, and
matrices are measurable functions of this retained history. Rollback and
marking are measurable operations, and the sigma-fields $\mathcal F_q^{\mathrm{ret}}$
continue to increase: answers observed in failed trials remain known.
Whenever an
operation installs an exact true value whose remaining digits have not
yet been queried, it first exposes and records that true suffix. This
also applies when the change is controlled by a deterministic estimate.

\begin{lemma}[Conditional freshness]\label{lem:conditional-freshness}
Suppose queries obey the preceding rule. Conditional on the retained
history, the unqueried true digits retain their product fair-sign law.
The same holds jointly for unused auxiliary signs. In particular, at
distinct sites whose remaining true suffixes are unqueried, those suffixes
are conditionally independent and have the uniform laws specified above.
After a revealed prefix of an independent uniform order, its remaining
order is conditionally uniform over the unexposed elements, jointly with
the stated product law for unqueried signs.

These statements hold after every deterministic finite number of
operations and at every almost surely finite stopping time of
$(\mathcal F_q^{\mathrm{ret}})$, with conditioning on the stopped sigma-field.
\end{lemma}

\begin{proof}
Condition on the history before a query. Its source index is then fixed.
Coordinates whose values are already recorded contribute no new
information. If a query includes both known and unqueried digits, apply
the following factorization to its unqueried part.
For a sign query, the product law factors into the queried coordinate and
the other unqueried coordinates. Adjoining the answer leaves the latter
factors unchanged. For a whole true suffix, apply the same argument to
its block of digits: test against finite cylinder functions of the other
coordinates and then use the monotone class theorem. This is a statement
about conditional expectations, so it also covers a continuous residual
answer without conditioning on an individual value of probability zero.
The same factorization applies to unused auxiliary signs.

For a uniform order with $r$ elements still unexposed, each possible
next element has exactly $(r-1)!$ completions. Its remaining order is
therefore uniform after that entry is revealed. An answer subsequently
queried at the selected index uses no unexposed order entry. Choices and
failure decisions made from the enlarged history add no information.
Induction proves the assertions after every fixed finite number of
operations.

For precision at a random operation time, let $\tau$ be a stopping time
of $(\mathcal F_q^{\mathrm{ret}})$ with $\PP(\tau<\infty)=1$, and let $\mathcal F_\tau^{\mathrm{ret}}$
be its stopped sigma-field. For $D\in\mathcal F_\tau^{\mathrm{ret}}$, apply each
finite-operation conditional-expectation identity against
$D\cap\{\tau=q\}\in\mathcal F_q^{\mathrm{ret}}$ and sum over $q$. These countably many
events exhaust $D$ up to a null set and give the same identity conditional
on $\mathcal F_\tau^{\mathrm{ret}}$. A whole-suffix query remains one operation in this
argument.
\end{proof}

\subsection{Schur reduction and the Ward Laplacian}

For the fixed $E$ and $\lambda>0$, almost surely every principal matrix
of the true $H_m-E$ is invertible. Indeed, for a nonempty subset its
determinant is a polynomial in the independent continuous potentials,
with nonzero highest monomial $\prod_{i}\lambda V_i$. Its zero set has
Lebesgue measure zero. There are finitely many subsets on the torus.
We henceforth work on this event, which also covers subsets chosen
adaptively later. This probability-one property is not a preliminary
test of unrevealed potentials and supplies no information to the
algorithm's filtration. The empty terminal block requires no inverse.

Let $\mathcal H$ be the real symmetric matrix with the current real
potential values on $A$ and exact true values on $T$. Eliminate the terminal
block at the real energy $E$, obtaining
\[
 M=(\mathcal H-E)_{AA}
   -(\mathcal H-E)_{AT}(\mathcal H-E)_{TT}^{-1}
                         (\mathcal H-E)_{TA}.
\]
The matrix $M$ is real symmetric. Its off-diagonal entries connect
nearest neighbors or sites joined through one connected component
of terminal sites. Throughout, the subscript in a Schur complement
names the \emph{eliminated} indices: if $R$ is the complement of $\mathcal E$,
\[
 \Schur_{\mathcal E} B
 =B_{RR}-B_{R\mathcal E}B_{\mathcal E\mathcal E}^{-1}B_{\mathcal ER}.
\]
We use this notation only when the displayed inverse exists.
Put $r=(r_i)_{i\in A}\in\CC^A$, $p_i=\Im r_i>0$,
and, for a scale $t>0$, define
\[
 g=(M-t\diag r)^{-1},\qquad h=\sqrt t\,g,\qquad
 K_{ij}=h_{ij}^2,\quad F_{ij}=|h_{ij}|^2,\quad q_i=\Im g_{ii}.
\]
The inverse exists because its matrix has strictly negative imaginary
part. In particular $\|g\|\le(t\min p_i)^{-1}$.
Both $g$ and $h$ are complex symmetric. For a matrix $B$,
$\Im B=(B-B^*)/(2i)$; here it is also the entrywise imaginary part.
The resolvent identity gives the Ward identity
\begin{equation}\label{eq:ward}
 h\diag(p)h^*=\Im g,\qquad Fp=q.
\end{equation}
Consequently
\[
 D=\diag(pq)-\diag(p)F\diag(p)
\]
is a real graph Laplacian, with nonnegative edge weights
$p_i p_jF_{ij}$. Its quadratic form is
\[
 D(f)=\frac12\sum_{i,j}p_i p_jF_{ij}|f_i-f_j|^2.
\]
The nonnegative weights $F_{ij}=|h_{ij}|^2$ describe how the Ward mass is
distributed across a row.
We use $B(f)=\langle f,Bf\rangle$ for the quadratic form of a Hermitian
matrix. At the committed states below, the kernel of $D$ is precisely
the constants. Its pseudoinverse $X=D^+$ is then an inverse on neutral
charges, namely vectors $y$ with $\sum_i y_i=0$.

We record the precise Schur identity for these inverse forms. If
$U$ is eliminated and $R=A\setminus U\ne\varnothing$, then $D_{UU}$
is positive definite: extending a zero-energy vector by zero on $R$
would otherwise give a nonconstant vector in $\ker D$. Moreover,
\[
 (\Schur_U D)(f_R)=\min_{f_U}D(f_R,f_U).
\]
It follows that $\Schur_U D$ has only the constant kernel. For a
neutral $y$ supported on $R$, extend it by zero to $A$ and use the
variational formula
\[
 X(y)=\sup_f\{2\Re\langle y,f\rangle-D(f)\}.
\]
Minimizing first over $f_U$ proves
\begin{equation}\label{eq:schur-inverse-form}
 (\Schur_U D)^+(y)=D^+(y).
\end{equation}
All later inverse-form comparisons across deletion are understood in
this sense, on survivor-supported neutral charges, not as an equality
of full pseudoinverse matrices.

At the beginning of a stage the state satisfies $r_i=g_{ii}$.
During a sweep or clearing trial, $r$ remains frozen. Thus $q=p$
at the stage boundary, and the difference $g_{ii}-r_i$ measures the
subsequent diagonal lag. Because $g$ itself depends on $r$, restoring
the equality requires a small self-consistent correction. The stage-end
reset in Section~\ref{sec:assembly} constructs that correction.

The role of $K$ is visible in this self-consistency equation. For the
physical active inverse, differentiation with respect to the complex
diagonal parameter gives
\[
 \frac{\partial g_{ii}}{\partial r_j}=t g_{ij}^2=K_{ij}.
\]
Thus the linearization of $r-(g_{ii}(r))_{i\in A}=0$ is $\Id-K$, which is
invertible whenever $\|K\|<1$.

For $n_0\le n<m$, a transition from scale $n$ to $n+1$ consists of an
ordinary sweep, any scheduled clearings at $t_{n+1}$, and a diagonal reset.
Clearings run only through $J=\lfloor m/2\rfloor$. At the end of stage $J$,
the construction aborts if any site remains pending; after $J$ it aborts
if the first trial at any tile fails. At scale $m$, a final residual sweep
inserts each remaining true suffix, leaves damping $\eta$, and has no reset.
Sections~\ref{sec:assembly} and~\ref{sec:history} specify the trial, marking,
and abort rules.

\subsection{The bounds carried between stages}

The full-torus positive nearest-neighbor Laplacian is denoted by
$\Lambda$. Thus $\Lambda(v)$ is the sum of $|v_i-v_j|^2$ over unordered
nearest-neighbor edges. Charges on a subset are extended by zero when
inserted in $\Lambda^+$. We fix
\[
 \delta=0.01,\qquad \kappa=0.005\delta,\qquad
 \nu=10^{-5}\delta/d,\qquad H_*=t^{1+\delta},\qquad b=t_{n_0}t.
\]
Thus $d\nu=10^{-5}\delta$; the dimension-scaled choice will absorb the
overlap of component neighborhoods. Constants in $O(\cdot)$ depend
only on $d$, the displayed fixed energy window, and the fixed broad
bounds specified below. By making $\lambda_0(d)$
smaller we may make $n_0$ arbitrarily large. Every later use of a strict
power margin is made after this single choice.

At a committed stage boundary, the entries of $r$ are bounded and
$c\le p_i\le C$. We maintain
\[
 \|K\|\le\tfrac13,\qquad
 \max_i\sum_j|h_{ij}|^4\le H_*,\qquad \ker D=\CC\mathbf1.
\]
The fourth-power bound is the squared row norm $\sum_jF_{ij}^2$.
It limits concentration across the columns of each row. The local
variance estimates need this control in addition to Ward's bound on
the weighted row sum.
On neutral charges we maintain the three comparisons
\begin{equation}\label{eq:comparators}
 \begin{aligned}
 X&\ge c_n b\,\Lambda^+,\\
 X&\le C_0\Id+C_n b\,\Lambda^+,\\
 X&\le(1+t^\kappa)\diag(p^{-2})+C_n t\,\Lambda^+.
 \end{aligned}
\end{equation}
Here $C_0$ is a sufficiently large fixed constant, while $c_n,C_n$
remain bounded above and below by positive constants, changing only
by summable relative errors. Boundedness of the weighted row sums of
$F$ also gives $X\ge c\Id$ on this subspace: the absolute row sums
of $D$ are at most $2\max_i p_iq_i$, so $D\le C\Id$.

The first two comparisons express an equivalence of spatial energies.
To formulate it also for the duplicated sites used in a sweep, allow
one or two \emph{ports} at each active physical site. For a port charge
$y$, let $y_\Sigma$ be the full-torus charge formed by summing its ports
at each physical site and extending by zero. A full-torus function $v$
is repeated at all ports at a site. Define the dual quadratic forms
\begin{equation}\label{eq:dual-forms}
 \begin{aligned}
 W^*(y)&=\|y\|_2^2+b\,\Lambda^+(y_\Sigma),\qquad \sum_x y_x=0,\\
 W(f)&=\inf_{f=\alpha+v|_{\mathrm{ports}}}
       \left(\|\alpha\|_2^2+b^{-1}\Lambda(v)\right).
 \end{aligned}
\end{equation}
The second form is on potentials modulo constants.
Lemma~\ref{lem:dual-forms} proves the duality and explains the resulting
relations $X\asymp W^*$ and $D\asymp W$. The final comparison in
\eqref{eq:comparators} is sharper: its high-frequency diagonal
coefficient tends to one. Retaining that coefficient will supply the
extra gain in the row-fourth-moment estimates.

We also track geometry. Let $d_T(i,j)$ be shortest-path pseudodistance in the
physical torus when an edge touching a terminal site has cost zero
and every other nearest-neighbor edge has cost one. The induction
will maintain
\[
 d_T(i,j)\ge\tfrac12\dist(i,j)-t^{-1.03}.
\]
The matrices $M$ can have large entries after a Schur reduction.
Nevertheless, every nonlocal entry joins points at zero $d_T$-distance.
This permits the resolvent decay estimate needed for local monitoring.
Section~\ref{sec:assembly} proves the geometric bound, including the
temporary terminal sites inside a clearing trial.

There are two distinct uses of these bounds. An ordinary trial stops at
the first local budget violation or failure of a required prefix bound,
including the global correction, row, and lag caps. The violating increment
is included in the stopped estimates.
Clearing prefixes are governed by their own stopped trial budgets, which
yield the provisional diagonal and row bounds; geometry is supplied
separately. For each
successful clearing, the assembly argument proves the $K$ and Laplacian
comparisons for the committed output from its stated hypotheses. The
trial statements make these two regimes explicit. The maximal
concentration estimate used for their stopped probability bounds is
proved in Section~\ref{sec:local-sweeps}.
